\documentclass[10pt]{article}
\usepackage[T1]{fontenc}
\usepackage[utf8]{inputenc}
\usepackage[french]{babel}
\usepackage{amsmath,amssymb,mathrsfs,mathpazo}
\usepackage{geometry}
\geometry{a4paper,margin=2cm}
\usepackage{graphicx,array,tabularx,tabularray,colortbl,multirow,multicol,booktabs,enumitem,pifont,ulem,fancybox,diagbox,cancel}
\usepackage{pstricks,pstricks-add,pst-plot,pst-tree,pst-node,pst-text,pst-func,pst-math,pst-bspline,pst-3dplot,pst-eucl,pst-solides3d}
\let\pscalculate\undefined
\usepackage{pst-bezier}
\usepackage{tikz,pgfplots,tkz-tab}
\usetikzlibrary{arrows,arrows.meta,patterns,calc,positioning,trees,decorations.markings,shapes.geometric,shapes.misc,shapes.arrows,intersections,decorations.pathmorphing,decorations.pathreplacing,backgrounds,shadows,babel}
\pgfplotsset{compat=1.18}
\usepgfplotslibrary{fillbetween}
\usepackage[np]{numprint}
\usepackage{listings,makecell,pas-tableur,twoopt}
\let\barre\undefined
\usepackage[pstricks]{tablvar}
\definecolor{darkgreen}{RGB}{0,100,0}
\providecommand{\red}{\color{red}}
\providecommand{\blue}{\color{blue}}
\providecommand{\green}{\color{green}}
\providecommand{\black}{\color{black}}
\providecommand{\white}{\color{white}}
\providecommand{\gray}{\color{gray}}
\providecommand{\cyan}{\color{cyan}}
\providecommand{\magenta}{\color{magenta}}
\providecommand{\yellow}{\color{yellow}}
\providecommand{\violet}{\color{violet}}
\providecommand{\orange}{\color{orange}}
\providecommand{\euro}{\texteuro}
\providecommand{\EUR}{\texteuro}
\providecommand{\ieme}{\textsuperscript{e}}
\providecommand{\degres}{\ensuremath{^\circ}}
\providecommand{\degree}{\ensuremath{^\circ}}
\providecommand{\vv}[1]{\overrightarrow{#1}}
\providecommand{\Coord}[1]{#1}
\providecommand{\decosix}{}
\providecommand{\starredbullet}{\ensuremath{\bullet}}
\pagestyle{empty}
\setlength{\parindent}{0pt}
\usepackage{adjustbox}\DeclareUnicodeCharacter{2236}{:}\DeclareUnicodeCharacter{2190}{\ensuremath{\leftarrow}}
\providecommand{\R}{\mathbb{R}}
\providecommand{\N}{\mathbb{N}}
\providecommand{\Z}{\mathbb{Z}}
\providecommand{\Q}{\mathbb{Q}}
\providecommand{\D}{\mathbb{D}}
\providecommand{\C}{\mathbb{C}}
\providecommand{\e}{{\mathrm{e}}}
\providecommand{\vect}[1]{\overrightarrow{#1}}
\providecommand{\vectt}[1]{\overrightarrow{\mathrm{#1}}}
\providecommand{\vvt}[1]{\overrightarrow{\mathrm{#1}}}
\providecommand{\np}[2][]{{#2~#1}}
\providecommand{\vv}[1]{{\overrightarrow{#1}}}
\renewcommand{\d}{{\mathrm{d}}}
\providecommand{\degres}{{^\circ}}
\providecommand{\degre}{{^\circ}}
\providecommand{\blue}{\color{blue}}
\providecommand{\red}{\color{red}}
\providecommand{\enskip}{\quad}
\providecommand{\llbracket}{\lbrack\!\lbrack}
\providecommand{\rrbracket}{\rbrack\!\rbrack}
\providecommand{\euro}{{\text{€}}}
\providecommand{\no}{{\text{n°}}}
\providecommand{\interval}[2]{{[#1;#2]}}
\providecommand{\intervFO}[2]{{[#1;#2[}}
\providecommand{\dotfill}{\ldots}

\providecommand{\ds}{\displaystyle}
\providecommand{\textding}[1]{#1}
\providecommand{\Oij}{\(\mathrm{O};\vec{\imath},\vec{\jmath}\)}
\providecommand{\Oijk}{\(\mathrm{O};\vec{\imath},\vec{\jmath},\vec{k}\)}
\providecommand{\Ouv}{\(\mathrm{O};\vec{u},\vec{v}\)}
\providecommand{\intervFF}[2]{\ensuremath{[#1;#2]}}
\renewcommand{\euro}{\ifmmode\text{\texteuro}\else\texteuro\fi}
\DeclareUnicodeCharacter{00B0}{\ifmmode^{\circ}\else\textdegree\fi}
\DeclareUnicodeCharacter{20AC}{\ifmmode\text{\texteuro}\else\texteuro\fi}
% Énoncés : Éducation nationale. Transcription et corrigés : APMEP, Denis Vergès et contributeurs. Découpage et classement : CoopMaths, Sébastien Lozano, Éric Elter et contributeurs. Adaptation de la présentation : Maths Explorer.
% CC BY-SA : https://coopmaths.fr/www/about/
% https://coopmaths.fr/alea/static/bac/2025/tex/bac_2025_06_sujet2_metropole_1_cor.tex
\begin{document}
\noindent\textit{APMEP / CoopMaths. CC BY-SA. Présentation adaptée par Maths Explorer.}\par\medskip
\medskip

\begin{enumerate}
\item Voici l'arbre des probabilités pondéré par les éléments de l'énoncé :

\begin{center}
\pstree[treemode=R,levelsep=3cm]{\TR{}}
{\pstree{\TR{$A~$}\taput{0,6}}
	{\TR{$B$} \taput{0,7}
	\TR{$\overline{B}$} \tbput{0,3}
	}
\pstree{\TR{$\overline{A}~$}\tbput{0,4}}
	{\TR{$B$} \taput{0,4}
	\TR{$\overline{B}$} \tbput{0,6}
	}
}
\end{center}
\item On sait que $P\left(\overline{A}\ \cap \overline{B}\ \right) = P\left(\overline{A}\right)  \times P_{\overline{A}}\left(\overline{B} \right) = 0,4 \times 0,6 = 0,24$.

La probabilité qu'une personne ne chute pas lors des deux premières séances est égale à 0,24.
\item $\left\{A~;~\overline{A}\right\}$ est un système complet d'évènements, donc d'après la formule des probabilités totales :

$P(B) = P(A \cap B) + P\left(\overline{A} \cap B\right) = P(A) \times P_A(B) + P\left(\overline{A}\right) \times P_{\overline{A}}(B)$, soit 

$P(B) = 0,6 \times 0,3 + 0,4 \times 0,4 = 0,18 + 0,16 = 0,34$.
\item D'après la formule de la probabilité conditionnelle :

$P_{\overline{B}}(\overline{A}) = \dfrac{P\left(\overline{A} \cap \overline{B}\right)}{P\left(\overline{B}\right)} = \dfrac{P\left(\overline{A} \cap \overline{B}\right)}{1 - P(B)} = \dfrac{0,24}{1 - 0,34} = \dfrac{0,24}{0,66} = \dfrac{24}{66} = \dfrac{4}{11}$ ou encore :

$P_{\overline{B}}\left(\overline{A}\right) \approx \np{0,3636}$, soit 0,364 au millième près.
\item 
	\begin{enumerate}
		\item On répète $100$ fois de façon identique et indépendante l'expérience de Bernoulli dont la probabilité de succès (ne pas avoir chuté lors des deux premières séances) est égale à 0,24.
		
La variable aléatoire $X$ suit donc la loi binomiale de paramètres $n = 100$ et $p = 0,24$.
		\item On a $P(X \geqslant 20) = 1 - P(X < 20) = 1 - P(X \leqslant 19)$ soit d'après la calculatrice 
		
$\approx 1 - 0,145$, d'où

		$P(X \geqslant 20) \approx 0,855$ au millième près.
		\item On sait qu'alors $E(X) = n \times p = 100 \times 0,24 = 24$.

Ceci signifie qu'en moyenne sur 100 personnes, 24 ne chuteront lors des deux premières séances.
		
	\end{enumerate}
\end{enumerate}

\medskip

\textbf{Partie B}

\begin{enumerate}
\item Puisque $T = T_1 + T_2$, par linéarité de l'espérance :

$E(T) = E(T_1 + T_2) = E(T_1) + E(T_2) = 40 + 60 = 100$.

Il faut donc attendre en moyenne 100 minutes sur les deux jours.

Les variables aléatoires $T_1$ et $T_2$ étant indépendantes on a de même :

$V(T) = V(T_1 + T_2) = V(T_1) + V(T_2)$

Comme $V(T_1) = \sigma(T_1)^2 = 10^2 = 100$ et que $V(T_2) = \sigma(T_2)^2 = 16^2 = 256$, on a donc 

$V(T) = 100 + 256 = 356$.
\item Soit $60 < T < 140$ soit en centrant sur l'espérance :

$60 - 100 < T - 100 < 140 - 100 \iff - 40 < T - 100 < 40 \iff |T - 100| < 40$ ou encore $|T - E(T)| < 40$.

D'après l'inégalité de Bienaimé-Tchebychev : 

$P(|T - E(T)| \geqslant 40) \leqslant \dfrac{V(T)}{40^2} \iff P(|T - E(T)| \geqslant 40) \leqslant \dfrac{V(T)}{40^2} \iff $

$P(|T - E(T)| < 40) \geqslant 1 - \dfrac{356}{\np{1600}} \iff P(|T - E(T)|  < 40) \geqslant \dfrac{\np{1244}}{\np{1600}}$.

Comme $\dfrac{\np{1244}}{\np{1600}} \approx \np{0,7775}$ qui est supérieur à 0,77, on a finalement :

\[P(60 < T < 140) > 0,77.\]

\end{enumerate}

\medskip

\end{document}
